% ============================================================
%  SECTION 1 — INTRODUCTION
% ============================================================
\section{Introduction}

Fractal geometry provides a mathematical framework for quantifying the structural complexity and statistical self-similarity of biological phenomena that classical Euclidean geometry fails to capture \cite{Mandelbrot1982}. In medical image analysis, the fractal dimension ($D$) has been widely adopted to characterize tissue morphology, demonstrating significant clinical utility in assessing pulmonary emphysema \cite{Mishima1999}, retinal vasculature \cite{Jelinek2000}, cortical folding \cite{Free1996}, and tumour margin irregularity \cite{Rangayyan2007}. Recently, fractal descriptors have evolved beyond standalone features, being integrated into hybrid machine learning pipelines to enhance model interpretability \cite{roberto2021fractal} and guide data augmentation strategies \cite{borgue2025association, Lopes2011}.

Despite this proven discriminative utility, existing literature focuses heavily on diagnostic effectiveness while largely neglecting the underlying computational trade-offs. Contemporary clinical datasets routinely comprise high-resolution images where algorithmic efficiency is no longer a secondary concern, but a strict practical constraint. Consequently, the critical research question must shift from \textit{whether} fractal dimension is discriminative to \textit{how efficiently and robustly} it can be extracted under realistic resource limitations.

This study investigates the relationship between theoretical algorithmic complexity, implementation choices, memory utilization, and practical execution time. We evaluate four fundamental estimation algorithms that span the principal design axes of fractal image analysis (binary versus grayscale inputs, and disjoint versus dense-overlapping covering strategies): Classical Box-Counting (BC) \cite{Mandelbrot1982}, Gliding Box (GB) \cite{Allain1991}, Differential Box-Counting (DBC) \cite{Sarkar1994}, and the Blanket Method (BM) \cite{Peleg1984}. 

We hypothesize that theoretical asymptotic improvements do not necessarily translate into practical runtime gains. While structural optimizations (such as integral images or sparse tables) reduce mathematical complexity, their real-world performance is heavily mediated by memory hierarchy bottlenecks, cache behavior, early-exit conditions, and preprocessing requirements. To test this hypothesis, this work addresses the following concise objectives:

\begin{itemize}
  \item \textbf{Theoretical vs. Empirical Scaling:} Derive formal asymptotic time and auxiliary space bounds for naive and structurally optimized formulations of all four methods, explicitly contrasting theoretical predictions with observed wall-clock performance on a standard medical dataset.
  \item \textbf{Resource Trade-offs:} Quantify the execution and memory overheads imposed by advanced data structures and threshold-based preprocessing steps, evaluating their feasibility for clinical deployment.
  \item \textbf{Discriminatory Effectiveness:} Assess the statistical separation between pathological and healthy tissues generated by each algorithm, culminating in a synthesized trade-off analysis that balances computational cost against diagnostic value.
\end{itemize}

The remainder of this article is structured as follows: Section~2 establishes the theoretical foundations of the evaluated methods. Section~3 details the algorithm designs and their asymptotic complexities. Section~4 outlines the experimental methodology. Section~5 analyzes the empirical results, explicitly addressing the divergence between theory and practice. Section~6 discusses threats to validity, and Section~7 concludes with evidence-driven guidelines for algorithm selection in medical image analysis.